% Held out of the report (2026-11-26, user): older DRTC results wait for the new robot results (E4).
% Not built; kept unchanged for later.

\chapter*{Abstract}
\addcontentsline{toc}{chapter}{Abstract}
\thispagestyle{plain}

Vision-language-action models are too large to run on the small computer that drives a low-cost robot arm, so they are
served from a remote GPU and every action depends on the network. When the network delays or loses messages, the
robot's queue of actions empties and the arm stops; the remote-inference client shipped with LeRobot~0.7 then ends the
run after one second without actions, and it never replaces a lost reply. This project studies that failure on an
SO-101 arm and tests whether Distributed Real-Time Chunking (\drtc), a scheduling method designed for unreliable
networks, recovers it, which of its parts matter, and how the motion can be kept smooth.

We built a platform on LeRobot~0.7 for recording, training, serving and evaluating policies across a laptop, a
Raspberry~Pi and two shared GPU servers; recorded and trimmed three datasets for two tasks; and trained seven policies
with their authors' recipes. We ported \drtc{} into LeRobot's remote-inference runtime with four fixes for slow links and
short chunks, implemented six smoothing methods and seven kinds of network fault, and built an evaluation framework with
frozen plans, seeded rounds, raw records of every trial, recovery and plan-deviation metrics, and primary comparisons
tested with the Holm correction. In a first robot run on a sorting task with \smolvla{} and \groot, a mix of a latency
spike, dropped replies and an outage lowered task success from \R{main/smolvla/baseline/M1} and
\R{main/groot/baseline/M1} to \R{main/smolvla/faults/M1}, the shipped client ending every faulted trial at the first
fault. With \drtc{} success was \R{main/smolvla/drtc/M1} and \R{main/groot/drtc/M1}, no trial was ended by a fault, and
the arm moved again within about a second and a half of every fault; smoothing reduced the jerk by factors of
\R{main/smolvla/jerkratio} and \R{main/groot/jerkratio} without lowering success. Because the shipped client gives up
by design, the experiments were redesigned to compare \drtc{} with a client that waits, to remove each of its parts in
turn, to choose a smoothing combination by a rule fixed in advance, and to cover two tasks and five policies on a
replay bench; the report gives their design and the results as they arrive.

\vspace{1em}
\noindent\textbf{Keywords:} robot learning, vision-language-action models, remote inference, action chunking,
network faults, real-time chunking, motion smoothing.
